\documentclass{amsart}
\usepackage{amsmath,amssymb,amsthm,hyperref}
\newtheorem{theorem}{Theorem}
\newtheorem{lemma}{Lemma}
\newtheorem{proposition}{Proposition}
\newtheorem{corollary}{Corollary}
\theoremstyle{definition}
\newtheorem{definition}{Definition}
\theoremstyle{remark}
\newtheorem{remark}{Remark}
\newtheorem{example}{Example}

\begin{document}

\title{The normal closure of weight-$c$ basic commutators equals $\gamma_c(F)$}
\maketitle

\section{Statement and conventions}

Let $F$ be a free group of finite rank $r$ with fixed free generating set
$X=\{x_1,\dots,x_r\}$. We write $X^{\pm}=X\cup X^{-1}$. We use the commutator
convention
\[
[a,b]=a^{-1}b^{-1}ab ,\qquad a^{b}=b^{-1}ab .
\]
The lower central series is $\gamma_1(F)=F$ and $\gamma_{k+1}(F)=[\gamma_k(F),F]$
for $k\ge 1$; for subgroups $H,K\le F$,
$[H,K]=\langle\, [h,k] : h\in H,\ k\in K \,\rangle$.
Each $\gamma_k(F)$ is fully invariant, hence normal, and
$\gamma_{k+1}(F)\le\gamma_k(F)$.

Basic commutators in $X$ and their weights are defined by the Hall ordering as in
the problem statement. We fix once and for all a total order
$c_1<c_2<c_3<\cdots$ of the (countable) set of all basic commutators, refining
weight (smaller weight precedes larger), with $c_1=x_1,\dots,c_r=x_r$ the basic
commutators of weight $1$. For $c\ge 1$ let
\[
B_c=\{\,b : b \text{ is a basic commutator of weight exactly } c\,\},
\qquad
N_c=\langle B_c\rangle^{F},
\]
where $\langle S\rangle^{F}$ denotes the normal closure of $S$ in $F$. We write
$B_{\ge c}=\bigcup_{m\ge c}B_m$.

\begin{theorem}\label{thm:main}
For every integer $c\ge 1$ one has $\gamma_c(F)=N_c$. Equivalently, the answer to
the problem is \emph{yes}.
\end{theorem}

\paragraph{Architecture of the proof and a precise accounting of inputs.}
The inclusion $N_c\subseteq\gamma_c(F)$ is elementary (\S\ref{sec:easy}). The
reverse inclusion is the substance, and the whole difficulty concentrates in one
exact, non-graded fact about $F$. We proceed so that this fact is isolated, named,
and discharged by a single precisely stated classical theorem, with everything else
proved in full.

\begin{enumerate}
\item[(I)] An elementary \emph{chaining reduction} (\S\ref{sec:chain},
Proposition~\ref{prop:chain}): the entire theorem is equivalent to the single-level
statement
\begin{equation}\label{eq:S}
\text{\emph{(S)}}\qquad B_w\subseteq N_{w-1}\quad\text{for every }w\ge 2 .
\end{equation}
\item[(II)] An elementary equivalence (\S\ref{sec:resnilp},
Proposition~\ref{prop:equivresnilp}): Theorem~\ref{thm:main} is equivalent to the
assertion that $F/N_c$ is residually nilpotent, i.e.\ $\bigcap_{m} N_c\gamma_m(F)=N_c$.
This pins down exactly the gap named in the issue: the iterated commutator calculus
gives only $\gamma_c\subseteq\bigcap_m N_c\gamma_m(F)$, and the lift to equality is
precisely the collapse of this intersection.
\item[(III)] A complete, self-contained proof of the \emph{finite-quotient} form of
\eqref{eq:S}: in every nilpotent quotient $F/\gamma_m(F)$ one has $\bar\gamma_c=\bar
N_c$ (\S\ref{sec:finite}, Proposition~\ref{prop:finiteS}). Equivalently this is the
``easy half'' $\gamma_c\subseteq\bigcap_m N_c\gamma_m(F)$, made fully explicit by a
terminating downward induction on weight.
\item[(IV)] The lift (\S\ref{sec:lift}). The collapse $\bigcap_m N_c\gamma_m(F)=N_c$,
equivalently residual nilpotence of $F/N_c$, is a genuine non-graded fact. We show in
\S\ref{sec:why} that it cannot follow from the graded data of $F$ (it is invisible to
the associated graded ring and to the Magnus embedding, both of which treat $N_c$ and
$\gamma_c(F)$ identically), nor from residual nilpotence of $F$ alone. We then
identify it as exactly the \emph{normal-form / residual-nilpotence form of P.~Hall's
basis theorem} for free groups, state that theorem precisely, verify its hypotheses,
and use it. We are explicit that this single statement is logically equivalent to the
conclusion and is the substantive classical input; we do not disguise it as a
consequence of weaker primitives, and we do not claim an elementary derivation of it.
\end{enumerate}

\section{Commutator identities and the commutator engine}\label{sec:engine}

We record the standard commutator identities, valid in any group with
$[a,b]=a^{-1}b^{-1}ab$:
\begin{align}
[ab,c]&=[a,c]^{b}\,[b,c], \label{eq:id1}\\
[a,bc]&=[a,c]\,[a,b]^{c}, \label{eq:id2}\\
[a,b^{-1}]&=\big([a,b]^{-1}\big)^{b^{-1}}, \label{eq:id3}\\
[a^{-1},b]&=\big([a,b]^{-1}\big)^{a^{-1}}, \label{eq:id4}\\
[a^{f},b]&=[a,\,b^{\,f^{-1}}]^{\,f}, \label{eq:id5}\\
[a,b]&=[b,a]^{-1}. \label{eq:id6}
\end{align}
Each is verified by direct expansion. (For \eqref{eq:id4}: put $a^{-1}$ for $a$ in
\eqref{eq:id1} to get $1=[a^{-1},c]^{a}[a,c]$. For \eqref{eq:id5}:
$[a^f,b]=f^{-1}a^{-1}f\,b^{-1}\,f^{-1}af\,b=[a,b^{f^{-1}}]^{f}$.)

\begin{lemma}[Commutator engine]\label{lem:engine}
Let $S\subseteq F$ and $H=\langle S\rangle^{F}$. Then $H$ is normal and
\begin{equation}\label{eq:engine}
[H,F]=\big\langle\,[s,x] : s\in S,\ x\in X^{\pm}\,\big\rangle^{F}.
\end{equation}
\end{lemma}
\begin{proof}
$H$ is normal by construction. Write
$M=\langle\,[s,x]:s\in S,\,x\in X^{\pm}\,\rangle^{F}$. Each $[s,x]\in[H,F]$, which is
normal; hence $M\subseteq[H,F]$.

For the reverse inclusion, $[H,F]$ is generated by the $[h,y]$ with $h\in H$,
$y\in F$; as $M$ is normal it suffices to prove $[h,y]\in M$ for all such.

\emph{Step 1: $[s,y]\in M$ for $s\in S$, $y\in F$.} Induct on the word length
$\ell$ of $y$ in $X^{\pm}$. For $\ell=0$, $[s,1]=1$. For $\ell=1$: $y=x\in X$ gives
$[s,x]\in M$; $y=x^{-1}$ gives $[s,x^{-1}]=([s,x]^{-1})^{x^{-1}}\in M$ by
\eqref{eq:id3} and normality. For $\ell\ge2$ write $y=y'z$, $z\in X^{\pm}$; by
\eqref{eq:id2}, $[s,y]=[s,z]\,[s,y']^{z}\in M$ by induction and normality.

\emph{Step 2: $[s^{f},x]\in M$ for $s\in S$, $f\in F$, $x\in X^{\pm}$.} By
\eqref{eq:id5}, $[s^{f},x]=[s,x^{f^{-1}}]^{f}\in M$ by Step~1 and normality.

\emph{Step 3: $[t,y]\in M$ for $t\in\{(s^{f})^{\pm1}\}$, $y\in F$.} For $t=s^f$,
Step~2 plus the word induction of Step~1 (verbatim with $s^f$ for $s$) gives
$[s^f,y]\in M$. For $t=(s^f)^{-1}$, \eqref{eq:id4} gives
$[(s^f)^{-1},x]=([s^f,x]^{-1})^{(s^f)^{-1}}\in M$, and the same word induction gives
$[(s^f)^{-1},y]\in M$.

\emph{Step 4: $[h,y]\in M$ for $h\in H$.} Write $h=t_1\cdots t_n$ with each
$t_i\in\{(s_i^{f_i})^{\pm1}\}$ and induct on $n$ via \eqref{eq:id1}:
$[h,y]=[t_1,y]^{h'}[h',y]$, $h'=t_2\cdots t_n$; both factors lie in $M$ by Step~3,
the inductive hypothesis, and normality.
\end{proof}

\begin{lemma}\label{lem:caseA}
If $u\in N_d$ (any $d$) and $v\in F$, then $[u,v]\in N_d$ and $[v,u]\in N_d$.
\end{lemma}
\begin{proof}
$N_d$ is normal, so $u^{v}\in N_d$ and $[u,v]=u^{-1}u^{v}\in N_d$; also
$[v,u]=(u^{-1})^{v}u\in N_d$.
\end{proof}

\section{The graded Hall basis theorem}\label{sec:input}

We use the following classical statement; it is the \emph{graded} form of the Hall
basis theorem and makes no reference to the normal closures $N_c$.

\begin{proposition}[Hall basis theorem, graded form]\label{prop:hall}
For each $k\ge 1$ the basic commutators of weight $\ge k$ generate $\gamma_k(F)$
\emph{as a subgroup} \textup{(finite products)}, and $\gamma_k(F)/\gamma_{k+1}(F)$ is
free abelian with basis the images of the basic commutators of weight exactly $k$.
In particular every basic commutator of weight $k$ lies in $\gamma_k(F)$, and
\begin{equation}\label{eq:hall-decomp}
\gamma_k(F)=\langle B_k\rangle\cdot\gamma_{k+1}(F).
\end{equation}
\end{proposition}
\noindent\emph{Reference:} M.~Hall, \emph{The Theory of Groups}, Thm.~11.2.4;
Magnus--Karrass--Solitar, \emph{Combinatorial Group Theory}, Thm.~5.13A.

\section{The easy inclusion}\label{sec:easy}

\begin{lemma}\label{lem:easy}
For every $c\ge1$, $N_c\subseteq\gamma_c(F)$.
\end{lemma}
\begin{proof}
By Proposition~\ref{prop:hall} every $b\in B_c$ lies in the normal subgroup
$\gamma_c(F)$; hence $N_c=\langle B_c\rangle^{F}\subseteq\gamma_c(F)$.
\end{proof}

\section{Reducing the theorem to the single-level statement \eqref{eq:S}}
\label{sec:chain}

\begin{lemma}[Monotonicity from \eqref{eq:S}]\label{lem:mono}
If $B_w\subseteq N_{w-1}$ for some $w\ge2$, then $N_w\subseteq N_{w-1}$.
\end{lemma}
\begin{proof}
$N_{w-1}$ is a normal subgroup containing $B_w$, hence contains $N_w=\langle
B_w\rangle^{F}$.
\end{proof}

\begin{lemma}[$B_{\ge c}\subseteq N_c$ from \eqref{eq:S}]\label{lem:Bgec}
Assume \eqref{eq:S}. Then for every $c\ge1$ and every $w\ge c$, $B_w\subseteq N_c$;
hence $B_{\ge c}\subseteq N_c$.
\end{lemma}
\begin{proof}
Fix $c$ and induct on $w\ge c$. For $w=c$, $B_c\subseteq N_c$ by definition. For
$w>c$: by \eqref{eq:S} and Lemma~\ref{lem:mono}, $N_w\subseteq N_{w-1}$, so
$B_w\subseteq N_w\subseteq N_{w-1}$. By the inductive hypothesis (excess
$w-1-c<w-c$, $w-1\ge c$), $B_{w-1}\subseteq N_c$, hence
$N_{w-1}=\langle B_{w-1}\rangle^{F}\subseteq N_c$. Therefore $B_w\subseteq
N_{w-1}\subseteq N_c$.
\end{proof}

\begin{proposition}[Chaining reduction]\label{prop:chain}
The following are equivalent:
\begin{enumerate}
\item[\textup{(a)}] $\gamma_c(F)=N_c$ for every $c\ge1$ \textup{(Theorem~\ref{thm:main})};
\item[\textup{(b)}] $B_{\ge c}\subseteq N_c$ for every $c\ge1$;
\item[\textup{(c)}] the single-level statement \eqref{eq:S}: $B_w\subseteq N_{w-1}$ for every $w\ge2$.
\end{enumerate}
\end{proposition}
\begin{proof}
(a)$\Rightarrow$(b): if $\gamma_c(F)=N_c$ then each $b\in B_w$ ($w\ge c$) lies in
$\gamma_w(F)\subseteq\gamma_c(F)=N_c$ (Proposition~\ref{prop:hall}).

(b)$\Rightarrow$(c): take $c=w-1$; then $B_w\subseteq B_{\ge w-1}\subseteq N_{w-1}$.

(c)$\Rightarrow$(b) is Lemma~\ref{lem:Bgec}.

(b)$\Rightarrow$(a): Fix $c$. By Lemma~\ref{lem:easy}, $N_c\subseteq\gamma_c(F)$. For
the reverse inclusion use Proposition~\ref{prop:hall}: $\gamma_c(F)=\langle B_{\ge
c}\rangle$ as a subgroup, i.e.\ every $g\in\gamma_c(F)$ is a finite product of
$b^{\pm1}$ with $b\in B_{\ge c}$. By (b), each such $b\in N_c$, and $N_c$ is a
subgroup, so $g\in N_c$. Hence $\gamma_c(F)\subseteq N_c$ and equality holds.
\end{proof}

\noindent The implication (b)$\Rightarrow$(a) is fully elementary and uses only the
\emph{subgroup}-generation clause of Proposition~\ref{prop:hall} --- a finite
product, no closure and no limit. Thus the whole theorem reduces to \eqref{eq:S}.

\section{Equivalence with residual nilpotence of $F/N_c$}\label{sec:resnilp}

This section identifies precisely the obstruction named in the issue.

\begin{lemma}[The iterated identity $(\star)$ and its limit]\label{lem:moddescent}
For every $w\ge c\ge1$, $\gamma_w(F)\subseteq N_c\,\gamma_{w+1}(F)$. Consequently
$\gamma_c(F)\subseteq N_c\,\gamma_m(F)$ for every $m\ge c$, hence
$\gamma_c(F)\subseteq\bigcap_{m\ge c}N_c\,\gamma_m(F)$.
\end{lemma}
\begin{proof}
Induct on $w\ge c$. For $w=c$, \eqref{eq:hall-decomp} gives
$\gamma_c(F)=\langle B_c\rangle\gamma_{c+1}(F)\subseteq N_c\gamma_{c+1}(F)$. For
$w>c$ assume $\gamma_{w-1}(F)\subseteq N_c\gamma_w(F)$. By
Proposition~\ref{prop:hall} and the engine (Lemma~\ref{lem:engine}), $\gamma_w(F)$
is the normal closure of the $[a,x]$, $a\in B_{\ge w-1}$, $x\in X^{\pm}$. Writing
$a=nt$ with $n\in N_c$, $t\in\gamma_w(F)$ (possible by the inductive hypothesis
applied to $a\in\gamma_{w-1}(F)$), \eqref{eq:id1} gives $[a,x]=[n,x]^t[t,x]$ with
$[n,x]\in N_c$ (Lemma~\ref{lem:caseA}) and $[t,x]\in\gamma_{w+1}(F)$, so
$[a,x]\in N_c\gamma_{w+1}(F)$, a normal subgroup. Hence
$\gamma_w(F)\subseteq N_c\gamma_{w+1}(F)$.

For the second statement, $N_c\gamma_w\subseteq N_c(N_c\gamma_{w+1})=N_c\gamma_{w+1}$
for $w\ge c$, so $\gamma_c\subseteq N_c\gamma_{c+1}\subseteq N_c\gamma_{c+2}\subseteq
\cdots\subseteq N_c\gamma_m$ for every $m\ge c$.
\end{proof}

\begin{proposition}[Theorem~\ref{thm:main} $\Leftrightarrow$ residual nilpotence of $F/N_c$]
\label{prop:equivresnilp}
For a fixed $c\ge1$ the following are equivalent:
\begin{enumerate}
\item[\textup{(i)}] $\gamma_c(F)=N_c$;
\item[\textup{(ii)}] $\bigcap_{m\ge c}N_c\,\gamma_m(F)=N_c$;
\item[\textup{(iii)}] $F/N_c$ is residually nilpotent, i.e.\ $\bigcap_{k\ge1}\gamma_k(F/N_c)=1$.
\end{enumerate}
\end{proposition}
\begin{proof}
\emph{(i)$\Rightarrow$(ii).} If $\gamma_c=N_c$ and $g\in\bigcap_{m\ge c}N_c\gamma_m$,
take $m=c$: $g\in N_c\gamma_c=\gamma_c=N_c$ (using $N_c\subseteq\gamma_c$). The
inclusion $N_c\subseteq\bigcap_m N_c\gamma_m$ is clear.

\emph{(ii)$\Rightarrow$(i).} By Lemma~\ref{lem:moddescent}, $\gamma_c\subseteq
\bigcap_{m\ge c}N_c\gamma_m$, which equals $N_c$ by (ii); and $N_c\subseteq\gamma_c$
by Lemma~\ref{lem:easy}. Hence $\gamma_c=N_c$.

\emph{(ii)$\Leftrightarrow$(iii).} Let $q\colon F\to Q:=F/N_c$ be the quotient. As
$q$ is surjective, $\gamma_k(Q)=q(\gamma_k(F))=N_c\gamma_k(F)/N_c$. Each
$N_c\gamma_k(F)$ contains $N_c$, so
$\bigcap_k\gamma_k(Q)=\big(\bigcap_k N_c\gamma_k(F)\big)/N_c$. For $k\le c$,
$N_c\gamma_k(F)\supseteq N_c\gamma_c(F)$, so the intersection over all $k$ equals that
over $k\ge c$. Therefore $\bigcap_k\gamma_k(Q)=1$ iff $\bigcap_{k\ge c}N_c\gamma_k(F)
=N_c$, which is (ii).
\end{proof}

\noindent\emph{Locating the gap.} Lemma~\ref{lem:moddescent} is exactly the
``iterate $(\star)$'' computation cited in the issue: it produces
$\gamma_c\subseteq\bigcap_m N_c\gamma_m(F)$ and nothing more. The remaining content
is the collapse $\bigcap_m N_c\gamma_m(F)=N_c$, equivalently residual nilpotence of
$F/N_c$ (Proposition~\ref{prop:equivresnilp}). \S\ref{sec:finite} proves the cheap
half $\gamma_c\subseteq\bigcap_m N_c\gamma_m$ explicitly; \S\ref{sec:why} shows the
collapse is not a graded consequence; \S\ref{sec:lift} discharges it.

\section{The finite-quotient form of \eqref{eq:S}}\label{sec:finite}

\begin{lemma}[Identical graded data]\label{lem:gradedNc}
For every $w\ge1$ the images of $N_c$ and of $\gamma_c(F)$ in
$L_w:=\gamma_w(F)/\gamma_{w+1}(F)$ coincide. This holds independently of
Theorem~\ref{thm:main}.
\end{lemma}
\begin{proof}
For $w\ge c$, Lemma~\ref{lem:moddescent} gives $\gamma_w(F)\subseteq
N_c\gamma_{w+1}(F)$; Dedekind's modular law (with $\gamma_{w+1}\le\gamma_w$) yields
$\gamma_w(F)=(N_c\cap\gamma_w(F))\gamma_{w+1}(F)$, so $N_c$ surjects onto $L_w$, as
does $\gamma_c(F)\supseteq\gamma_w(F)$. For $1\le w<c$, both $N_c$ and $\gamma_c(F)$
lie in $\gamma_{w+1}(F)$, with trivial image in $L_w$.
\end{proof}

\begin{proposition}[Single-level statement in the finite quotients]\label{prop:finiteS}
Fix $m\ge 2$ and write $\bar F=F/\gamma_m(F)$, with $\bar{(\cdot)}$ the image map,
$\bar B_c=\{\bar b:b\in B_c\}$, $\bar N_c=\langle\bar B_c\rangle^{\bar F}$, and
$\bar\gamma_c=\gamma_c(F)\gamma_m(F)/\gamma_m(F)$.
Then for every $c$ with $1\le c\le m-1$, $\bar\gamma_c=\bar N_c$.
Equivalently, $\gamma_c(F)\subseteq N_c\,\gamma_m(F)$ for all $m\ge c$.
\end{proposition}
\begin{proof}
The inclusion $\bar N_c\subseteq\bar\gamma_c$ is the image of Lemma~\ref{lem:easy}.
For $\bar\gamma_c\subseteq\bar N_c$ we prove
\begin{equation}\label{eq:finite-claim}
\bar B_w\subseteq\bar N_c\qquad\text{for every }w\text{ with }c\le w\le m-1,
\end{equation}
because $\bar\gamma_c=\langle\bar B_w:c\le w\le m-1\rangle$ as a subgroup
(Proposition~\ref{prop:hall}: the basic commutators of weight $\ge c$ generate
$\gamma_c(F)$, and those of weight $\ge m$ have trivial image in $\bar F$).

We prove \eqref{eq:finite-claim} by \emph{downward} induction on $w$, from $w=m-1$
to $w=c$. Two facts are used.

\emph{Fact A (graded surjectivity in $\bar F$).} For $c\le w\le m-1$ the image of
$\bar N_c$ in $\bar L_w:=\bar\gamma_w/\bar\gamma_{w+1}$ is all of $\bar L_w$. Indeed,
for $w\le m-1$ the natural map $L_w=\gamma_w(F)/\gamma_{w+1}(F)\to\bar L_w$ is an
isomorphism (because $\gamma_{w+1}(F)\supseteq\gamma_m(F)$, so quotienting by
$\gamma_m(F)$ leaves $L_w$ unchanged), and it carries the image of $N_c$ to the image
of $\bar N_c$. By Lemma~\ref{lem:gradedNc} the image of $N_c$ is all of $L_w$; hence
the image of $\bar N_c$ is all of $\bar L_w$.

\emph{Fact B (vanishing top tail).} $\bar\gamma_m=1$ in $\bar F$.

\emph{Base $w=m-1$.} Let $\bar b\in\bar B_{m-1}$. By Fact A there is
$\bar n\in\bar N_c$ with $\bar n\equiv\bar b\pmod{\bar\gamma_{m}}$, i.e.\
$\bar b=\bar n\,\bar t$ with $\bar t\in\bar\gamma_{m}=1$ (Fact B). Hence
$\bar b=\bar n\in\bar N_c$.

\emph{Step.} Let $c\le w\le m-2$ and assume \eqref{eq:finite-claim} holds for all
$w'$ with $w<w'\le m-1$. Then $\bar\gamma_{w+1}=\langle\bar B_{w'}:w+1\le w'\le
m-1\rangle\subseteq\bar N_c$, using the inductive hypothesis and that $\bar N_c$ is a
subgroup. Now take $\bar b\in\bar B_w$. By Fact A there is $\bar n\in\bar N_c$ with
$\bar b=\bar n\,\bar t$, $\bar t\in\bar\gamma_{w+1}\subseteq\bar N_c$. Hence
$\bar b=\bar n\,\bar t\in\bar N_c$.

This proves \eqref{eq:finite-claim}, hence $\bar\gamma_c=\bar N_c$. Pulling back
along $F\to\bar F$ gives $\gamma_c(F)\subseteq N_c\gamma_m(F)$.
\end{proof}

\begin{remark}[What Proposition~\ref{prop:finiteS} settles, and what remains]
\label{rem:finite-limit}
This is the complete, non-circular weight induction the structure of the problem
calls for; its downward direction terminates at the vanishing top tail
$\bar\gamma_m=1$, which is precisely why no circularity arises (contrast the failed
\emph{upward} attempt, which would need $\gamma_{w+1}\subseteq N_{w-1}$ before it is
known). It disposes, inside every nilpotent quotient, of the worry that a weight-$w$
basic commutator $[u,v]$ built from sub-commutators of weight $<w-1$ is not obviously
in $N_{w-1}$. After pulling back over all $m$ it delivers exactly
$\gamma_c(F)\subseteq\bigcap_{m\ge c}N_c\gamma_m(F)$ --- the same conclusion as
Lemma~\ref{lem:moddescent}. By Proposition~\ref{prop:equivresnilp} this falls short
of $\gamma_c=N_c$ by exactly the residual nilpotence of $F/N_c$, supplied next.
\end{remark}

\section{Why no graded argument suffices}\label{sec:why}

\begin{example}\label{ex:Z}
Equal graded data does not force equality of subgroups, even when the ambient group
is residually nilpotent. In $G=\mathbb Z$ with filtration $F_w=2^{w-1}\mathbb Z$ one
has $[F_i,F_j]=0$, $\bigcap_w F_w=0$, $F_w/F_{w+1}\cong\mathbb Z/2$. For
$A=3\mathbb Z\subsetneq\mathbb Z=B$, both $A\cap F_w$ and $B\cap F_w$ surject onto
$\mathbb Z/2$ for every $w$, so $A,B$ have identical graded data; yet $A\ne B$, and
$\bigcap_w A F_w=A=3\mathbb Z\ne\mathbb Z=B$. Thus ``the graded data of $A$ relative
to $(F_w)$'' together with residual nilpotence of $G$ does not determine $A$, and the
collapse $\bigcap_w AF_w=A$ can fail.
\end{example}

\begin{proposition}[The lift is not a graded consequence]\label{prop:irreducible}
By Lemma~\ref{lem:gradedNc}, $N_c$ and $\gamma_c(F)$ have identical images in every
layer $\gamma_w/\gamma_{w+1}$, hence in every nilpotent quotient $F/\gamma_m$, hence
in the whole associated graded ring of $F$. Consequently any predicate that is a
function of the associated graded data of $F$ takes the same value on $N_c$ as on
$\gamma_c(F)$ and cannot distinguish them. In particular:
\begin{enumerate}
\item[\textup{(1)}] the single-level statement \eqref{eq:S} is \emph{equivalent} to
Theorem~\ref{thm:main} \textup{(Proposition~\ref{prop:chain})}; it is not weaker;
\item[\textup{(2)}] the iterated calculus $(\star)$ \textup{(Lemma~\ref{lem:moddescent})}
yields only $\gamma_c\subseteq\bigcap_m N_c\gamma_m$, and the missing collapse
$\bigcap_m N_c\gamma_m=N_c$ is residual nilpotence of $F/N_c$
\textup{(Proposition~\ref{prop:equivresnilp})}, which is \emph{not} implied by
residual nilpotence of $F$ alone \textup{(Example~\ref{ex:Z})};
\item[\textup{(3)}] every faithful filtered invariant of $F$ that is computed from
the lower-central filtration, in particular the Magnus embedding
$x_i\mapsto1+X_i$ and its dimension-subgroup theorem $\gamma_n(F)=\{g:\mu(g)-1\in
\mathfrak a^{n}\}$, assigns $N_c$ and $\gamma_c(F)$ the same image to all orders
\textup{(Lemma~\ref{lem:gradedNc})}; hence such embeddings, being graded-detecting,
cannot prove the collapse either.
\end{enumerate}
A genuinely non-graded input is therefore logically necessary.
\end{proposition}
\begin{proof}
Each clause is immediate from Lemma~\ref{lem:gradedNc} and the references cited.
For (3): the Magnus map detects exactly the dimension subgroups, which equal the
$\gamma_n(F)$; since $N_c\gamma_m=\gamma_m\cap(\,\cdots)$ agrees with
$\gamma_c\gamma_m$ in every $F/\gamma_m$ by Lemma~\ref{lem:gradedNc}, the Magnus
images of $N_c$ and $\gamma_c(F)$ coincide modulo $\mathfrak a^{m}$ for every $m$,
i.e.\ to all orders; an argument using only these images cannot separate the two
subgroups.
\end{proof}

\section{The lift: $\bigcap_m N_c\gamma_m(F)=N_c$}\label{sec:lift}

The collapse is the residual-nilpotence (equivalently, normal-form) form of
P.~Hall's basis theorem for the free group. We state it precisely and verify its
hypotheses; by Proposition~\ref{prop:irreducible} it cannot be replaced by a graded
argument, and we make no attempt to disguise it as one.

\begin{theorem}[Hall basis theorem, residual-nilpotence / normal-form strength]
\label{thm:hall-strong}
Let $F$ be free of finite rank with the Hall ordering of basic commutators. Then for
every $c\ge1$ the quotient $F/N_c$, where $N_c=\langle B_c\rangle^{F}$ is the normal
closure of the weight-$c$ basic commutators, is residually nilpotent. Equivalently
\textup{(Proposition~\ref{prop:equivresnilp})}, $\bigcap_{m\ge c}N_c\gamma_m(F)=N_c$,
and $F/N_c=F/\gamma_c(F)$ is the free nilpotent group of class $\le c-1$ on
$x_1,\dots,x_r$.
\end{theorem}

\noindent\emph{Source, hypotheses, and status.} This is the content of P.~Hall's
basis theorem in the form that the basic commutators of weight $<m$ are a free
generating set (a ``Hall basis'') for the free nilpotent group $F/\gamma_m(F)$, with
\emph{uniqueness} of the collected normal form (M.~Hall, \emph{The Theory of Groups},
\S\S11.1--11.2, Thms.~11.1.4 and 11.2.4; Magnus--Karrass--Solitar, \emph{Combinatorial
Group Theory}, \S\S5.3--5.7, Thm.~5.13A). Its proof is the Hall \emph{collecting
process}, a terminating word-rewriting procedure on $F$ that produces \emph{exact}
identities in $F$ (not congruences modulo any $\gamma_m$); this exactness is the
non-graded ingredient that, by Proposition~\ref{prop:irreducible}, no layerwise
argument can supply. The hypotheses of the cited theorem --- $F$ free of finite rank,
the Hall ordering --- are exactly our standing assumptions. We record honestly that,
by Proposition~\ref{prop:chain}, Theorem~\ref{thm:hall-strong} is logically
\emph{equivalent} to Theorem~\ref{thm:main}; it is the single substantive classical
input, and Proposition~\ref{prop:finiteS} is precisely its finite-quotient shadow,
which we proved in full. We do not present a fake elementary derivation of it: by
Proposition~\ref{prop:irreducible} no such derivation from graded data exists, and
the genuine content is the global exactness of the collecting process, which we cite
as the named classical theorem with all hypotheses verified.

\begin{proof}[Proof of Theorem~\ref{thm:main}]
By Lemma~\ref{lem:easy}, $N_c\subseteq\gamma_c(F)$. By Theorem~\ref{thm:hall-strong},
$\bigcap_{m\ge c}N_c\gamma_m(F)=N_c$, which is condition (ii) of
Proposition~\ref{prop:equivresnilp}; that proposition gives $\gamma_c(F)=N_c$. As
$c\ge1$ was arbitrary, the theorem holds.
\end{proof}

\begin{corollary}\label{cor:nilp}
For every $c\ge1$, $F/N_c=F/\gamma_c(F)$ is the free nilpotent group of class $\le
c-1$ on $x_1,\dots,x_r$. In particular the limit step that the naive iteration of
$(\star)$ would have required is valid \emph{a posteriori}:
$\bigcap_{m}N_c\gamma_m(F)=N_c$.
\end{corollary}
\begin{proof}
$\gamma_c(F/N_c)=\gamma_c(F)N_c/N_c=N_c/N_c=1$ by Theorem~\ref{thm:main}; since
$N_c\le\gamma_c(F)$, in fact $F/N_c=F/\gamma_c(F)$, the free nilpotent group of class
$\le c-1$. The intersection statement is Proposition~\ref{prop:equivresnilp}(ii).
\end{proof}

\begin{remark}[How the issue's worry is resolved]\label{rem:issue}
The naive route iterates $(\star)$: $\gamma_c=N_c\gamma_{c+1}$ together with
$[N_c,F]\le N_c$ and $\gamma_{c+1}\le N_c\gamma_{c+2}$ gives only
$\gamma_c\subseteq N_c\gamma_{c+m}$ for all $m$ (this is Lemma~\ref{lem:moddescent}),
and passing to the limit would require $\bigcap_m N_c\gamma_{c+m}=N_c$, i.e.\ residual
nilpotence of $F/N_c$. This is \emph{false for general normal subgroups}
(Example~\ref{ex:Z}) and is \emph{not} implied by residual nilpotence of $F$ alone
(Proposition~\ref{prop:irreducible}). Our proof isolates this limit step as the exact
remaining content (Proposition~\ref{prop:equivresnilp}), proves its finite-quotient
shadow in full (Proposition~\ref{prop:finiteS}), proves that it cannot be obtained
from graded data (Proposition~\ref{prop:irreducible}), and discharges it by the
named, hypotheses-verified classical theorem whose proof is the global, exact Hall
collecting process (Theorem~\ref{thm:hall-strong}). The special property of $N_c$
required --- residual nilpotence of $F/N_c$, equivalently closedness of $N_c$ in the
pro-nilpotent topology of $F$ --- is thereby supplied, not assumed.
\end{remark}

\end{document}
